% Section 5: Applications
\subsection{Applications}
\label{sec:forecasting_rl_applications}

We outline four application domains where the prediction-and-decision interaction
is the active research question. The unifying claim across them is that joint or
decision-aware estimation outperforms two-step pipelines whenever the forecast
distribution is misspecified or the decision is asymmetric, and the unifying
methodological move is to identify the structural reason why the integrated
approach should work and to put a finite-sample or asymptotic guarantee under
it.

\subsubsection{Inventory and the contextual newsvendor}
\label{sec:app_inventory}

The contextual newsvendor is the canonical decision-under-uncertainty problem
and the cleanest empirical case for decision-focused learning. Given an
unknown future demand $D$, an order quantity $q$, and asymmetric per-unit
costs $b$ for backorders and $h$ for holding, the realised cost
$C(q, D) = b (D - q)^+ + h (q - D)^+$ is minimised at the
$b / (b + h)$ quantile of the demand distribution. \citet{BanRudin2019}
re-cast the problem when the decision-maker has access to a high-dimensional
feature vector $x \in \mathbb R^p$ that may include weather, seasonality,
economic indicators, and customer demographics. They propose two single-step
data-driven estimators that compute the order quantity directly from
$(x, d)$-pairs without an intermediate demand-distribution estimation step.
The Empirical Risk Minimisation procedure with a linear decision rule
$q(x) = \beta^\top x$ solves
\begin{equation}
\hat\beta_n = \arg\min_{\beta \in \mathbb R^p}
  \frac{1}{n} \sum_{i=1}^n
    \bigl[ b (d_i - \beta^\top x_i)^+ + h (\beta^\top x_i - d_i)^+ \bigr]
  + \lambda\, \Omega(\beta),
\label{eq:banrudin_erm}
\end{equation}
where $\Omega$ is an $\ell_0$, $\ell_1$, or $\ell_2$ regulariser yielding a
mixed-integer, linear, or quadratic programme respectively. The Kernel-weights
Optimisation procedure uses Nadaraya-Watson weights $w_i(x) \propto K((x_i - x)
/ h)$ and returns the empirical $b / (b + h)$ quantile of the weighted demand
distribution via a simple sorting algorithm.

\citet{BanRudin2019} prove finite-sample performance bounds on the
out-of-sample cost of both estimators. Their Theorem~4 establishes that the
Sample Average Approximation procedure without features has a constant
$O(1)$ bias that does not shrink in $n$, while the feature-based ERM
procedure attains a regret that scales with $\sqrt{p / n}$, which generalises
the no-feature bound of \citet{Levi2007} to the high-dimensional contextual
case. The empirical case for the integrated approach is on UK hospital nurse
staffing data, where the demand for emergency-room beds depends on
day-of-week, seasonal, and weather features. The best ERM result, with
$\ell_1$ regularisation, attains a cost improvement of $23\%$ relative to the
best-practice day-of-week-clustered SAA benchmark, equivalent to a saving of
$\pounds 44{,}219$ per ward per year. The best KO result attains $24\%$ and
$\pounds 46{,}555$ per ward per year, both statistically significant at the
$5\%$ level, and the KO sorting procedure runs in $0.05$ seconds versus the
two-orders-of-magnitude slower SAA baseline. The chapter's thesis quote is
from the introduction. ``Rather than a two-step process of first estimating
a demand distribution then optimizing for the optimal order quantity, we
propose solving the Big Data newsvendor problem via single-step machine
learning algorithms.''

\citet{Oroojlooyjadid2020} extend the \citet{BanRudin2019} programme from
linear to neural-network decision rules. The deep feed-forward network
$\theta(x)$ takes the feature vector as input and outputs the order quantity
directly, with the loss function being the newsvendor cost
$\mathcal L(\theta) = \sum_i [c_u (d_i - \theta(x_i))^+ + c_o
(\theta(x_i) - d_i)^+]$ rather than a forecast-error loss. The deep network
generalises the linear hypothesis class of \citet{BanRudin2019} and is
particularly effective for problems with small data, unknown demand
distributions, or volatile historical data. The taxonomic contribution of
\citet{Oroojlooyjadid2020} is the five-approach typology that classifies the
literature on the feature-based newsvendor, ranging from estimate-as-solution
through parametric estimation plus critical fractile and weighted-SAA via
machine learning, to the linear ML programme of \citet{BanRudin2019} and the
deep-learning generalisation of the present paper. The taxonomy makes the
chapter-level argument concrete. The deepest forecasting-for-decision
approaches are the ones that train the predictor on the realised decision
cost, and the empirical advantage over predict-then-optimise grows with the
level of demand-distribution misspecification.

The deep-generative extension of \citet{HanZhengMisra2024} reformulates the
contextual newsvendor under price-dependent demand and joint pricing-and-inventory
decisions. They observe that when price $p$ varies in the historical data,
traditional ERM is infeasible because the demand distribution is
decision-dependent, namely $D \mid X = x, P = p$ varies with the price.
The remedy is to estimate the conditional density $f_{D \mid X, P}$ by a
conditional deep generative model trained on the empirical
$\{(x_i, p_i, d_i)\}$ tuples, and then to sample from the estimated density
to construct an empirical conditional profit function
$\hat\pi(x, p, q) = \mathbb E_{D \sim \hat f}[(p - c) \min(D, q) - h (q - D)^+]$.
Their three consistency theorems establish that the cDGM-based inventory
decision converges to the true conditional optimum for any price, that the
estimated profit function is consistent for the conditional expected profit,
and that the joint pricing-and-inventory decisions are asymptotically
optimal under the assumption that the conditional generator converges to
the true conditional density. The empirical validation on Kaggle food
demand data confirms the consistency results and shows that the
generative approach dominates the linear and tree-ensemble baselines on
realised profit. The methodological message is that decision-dependent
uncertainty, which arises whenever the decision affects the distribution of
the outcome, cannot be addressed by a fixed-price ERM and requires the
joint distributional modelling that a conditional generative model provides.

\citet{Amazon2025}, working at Amazon Supply Chain Optimization
Technologies, push the framework to the multi-echelon production setting and
treat the forecaster as a differentiable simulator inside an RL inventory
ordering loop. Their Drain Model emulates the joint distribution
$p(o^i_t, c^i_t \mid H^i_t, \theta)$ of outbound units $o^i_t$ and shipping
costs $c^i_t$ for product $i$ at time $t$, conditional on the historical
state $H^i_t$ that includes inventory levels, customer glance views, active
warehouse set, and exogenous covariates. The output distribution combines a
discrete part with explicit probability mass on small integer outbound values
and a quantile-based tail for large outbounds, and the model is trained on a
linear combination of cross-entropy, quantile, and negative-log-likelihood
losses. The negative-log-likelihood term is preferred over the quantile loss
for the trajectory-generation case because the log-likelihood is a strictly
proper scoring rule in the sense of \citet{GneitingRaftery2007} while the
quantile loss is only proper for a single marginal prediction. The core
production challenge is that historical data only captures outcomes under the
current inventory placement policy, while a learned RL ordering policy will
visit out-of-distribution inventory configurations. \citet{Amazon2025} address
this with an off-policy backtesting oracle that uses conditional glance-view
conversion models to generate counterfactual fulfilment scenarios. The
authors note that ``preliminary results demonstrate the model's accuracy
within the in-distribution setting,'' which is the closest the published
literature gets to a TSFM-style forecaster used as a differentiable simulator
for an inventory RL agent, and the explicit acknowledgement of the
off-policy out-of-distribution problem is the operational analogue of the
Lucas critique discussed in Section~\ref{sec:bridge}.

\subsubsection{Electricity dispatch and microgrid operation}
\label{sec:app_electricity}

Electricity dispatch is the cleanest empirical decision-focused win in the
applied literature. The downstream optimisation is a two-stage robust
operation problem with binary unit-commitment variables, energy-storage
state-of-charge constraints, ramping limits, and demand-side flexibility, and
the forecast objects are the joint distributions of renewable generation,
load demand, and electricity price. Two-stage and stochastic dispatch
problems with a downstream linear or convex programming structure are exactly
the setting for which the SPO and SPO$^+$ machinery of
Section~\ref{sec:dfl} provides finite-sample guarantees, and the empirical
record bears out the theoretical promise.

\citet{DontiAmosKolter2017} provide the first task-based-learning demonstration
on real grid data, namely the PJM Interconnection electrical load series at
hourly resolution over seven years of training data and a year and three quarters
of testing data. The generator-scheduling problem has a $24$-hour horizon
with explicit ramping constraints and asymmetric over- and under-generation
penalties. A two-hidden-layer feed-forward neural network predicts the
hourly mean and variance of next-day load under a Gaussian assumption, and the
KKT-differentiable convex relaxation of the scheduling problem provides the
gradient pathway through which the task loss propagates to the network
weights. The task-based net reduces realised scheduling cost by $38.6\%$
relative to the RMSE-trained net using the same optimisation routine, and by
$8.6\%$ relative to a cost-weighted-RMSE variant. The result establishes that
decision-focused training improves realised operational cost even when the
predictor is structurally well-specified (Gaussian load with constant
variance) and only the variance heteroscedasticity is misspecified.

\citet{CaoXu2025} extend the decision-focused programme to multi-stage
microgrid operation by embedding a multilayer encoder-decoder probabilistic
forecaster inside a two-stage robust optimisation solved by the
column-and-constraint generation algorithm. The forecaster handles the
joint distribution of load demand, photovoltaic power, and wind turbine
power at $15$-minute resolution. The architecture is hierarchical, with
residual blocks followed by parallel encoder-decoder chains, Swish
activations, and a temporal attention decoder that projects encoded features
back into the prediction horizon. The training loss combines the continuous
ranked probability score with a regret-type decision-focused term
\begin{equation}
\mathrm{Regret}_{\mathrm{SPO}}(D, \hat D)
= J\bigl(\bar z_\rho(\hat D), D\bigr) - J\bigl(z^\star(D), D\bigr),
\label{eq:caoxu_regret}
\end{equation}
wrapped in a softplus operator for numerical stability. The downstream
operation problem includes binary commitment variables and is therefore
non-differentiable, so \citet{CaoXu2025} introduce a continuous convex
surrogate by relaxing the binaries to $[0, 1]$ and adding a Tikhonov
regularisation $\rho \|z\|^2$ to guarantee strong convexity. Gradients
through the surrogate optimiser $\bar z_\rho(\hat D)$ are obtained by
implicit differentiation of the Karush-Kuhn-Tucker residual, in the spirit
of \citet{DontiAmosKolter2017}. The forward decision evaluation uses the
exact mixed-integer TSRO via C\&CG, and only the gradient path runs through
the convex surrogate.

The empirical results of \citet{CaoXu2025} are on modified IEEE 33-bus and
69-bus power distribution networks with two ESS units and three to four
PV and wind generators per network. On the forecasting side, the MED model
attains $R^2 = 0.9651$ and $\mathrm{CRPS} = 0.72$ on load demand, $R^2 =
0.9843$ on PV power, and $R^2 = 0.9852$ on wind power, all dominating
DeepAR, Transformer, and Bi-LSTM baselines on all four metrics across the
three series. On the dispatch side, the breakdown of daily operation cost in
Table IV of \citet{CaoXu2025} shows the chapter's punchline cleanly. The
baseline two-stage forecast-then-optimise pipeline costs \$5{,}127 per day on
the 33-bus system. Replacing the deterministic forecast with DeepAR shaves
$5.9\%$ off the net cost; adding the MED probabilistic forecaster shaves
$11.5\%$; and switching to the end-to-end MED-plus-SPO-plus-TSRO framework
brings the reduction to $18.0\%$. On the larger 69-bus system the
corresponding reductions are $4.6\%$, $8.4\%$, and $12.6\%$ from a baseline of
\$7{,}521 per day. The forecast-only improvements (DeepAR and MED without
SPO) account for roughly two thirds of the total reduction; the
decision-focused training contributes the remaining third. The robustness
analysis varies the uncertainty budget $\Gamma$ and shows that the
end-to-end framework's $95\%$ confidence-interval width on total operation
cost is $42\%$ narrower than the baseline's, which is the operational
analogue of the calibration improvement that the SPO programme is designed
to deliver.

The chapter's thesis statement is the abstract sentence of \citet{CaoXu2025}.
The proposed method aligns forecasting accuracy with operational objectives
by directly minimising decision regret via a surrogate SPO loss function. The
combination of the formal regret bound from \citet{LiuGrigas2021},
the implicit-KKT-differentiation routine from \citet{DontiAmosKolter2017}, and
the empirical $12$-$18\%$ cost reduction on operational benchmarks makes
electricity dispatch the cleanest case study in the chapter for the practical
value of decision-focused forecasting.

\subsubsection{Macroeconomic nowcasting and monetary policy reaction}
\label{sec:app_macro}

Macroeconomic nowcasting and monetary policy reaction sit at the intersection
of the chapter's two halves. The forecasting half is anchored in
Section~\ref{sec:adp_forecasting} by \citet{Fernandes2024}, who establishes
that a neuro-dynamic-programming nowcaster trained by TD(0) on a 26-country EU
panel transfers to a held-out country and lowers out-of-sample MAE by $32\%$
relative to OLS on Portuguese GDP over 2015Q1-2023Q4. The decision half is
anchored by \citet{HinterlangTaenzer2024}, who use the Bundesbank Discussion
Paper 51/2021 to compute optimal interest-rate reaction functions by
reinforcement learning under realistic central-bank loss functions, including
the zero lower bound and nonlinear preferences. Their setup is the cleanest
existing template for an RL agent in a calibrated macroeconomic environment.

The environment is a two-equation system in inflation $\pi_t$ and the output
gap $y_t$, evolving according to estimated transition equations
\begin{equation}
y_t = f_y(s^y_t) + \varepsilon^y_t, \qquad
\pi_t = f_\pi(s^\pi_t) + \varepsilon^\pi_t,
\label{eq:hinterlang_env}
\end{equation}
where the state vectors $s^y_t$ and $s^\pi_t$ contain contemporaneous and
lagged values of $y, \pi$, and the nominal interest rate $i_t$. The state
forms a New-Keynesian-three-equation representation in the spirit of
\citet{RotembergWoodford1997}, with the central-bank reaction function
explicitly excluded from $f_y$ and $f_\pi$. The functional forms $f_y, f_\pi$
are estimated in two variants. The linear variant is a restricted structural
vector autoregression of order two, with the restrictions chosen so that the
Bayes and Akaike information criteria favour the parsimonious specification
over the unrestricted SVAR(2). The nonlinear variant approximates each
transition equation by a single-hidden-layer feed-forward neural network with
the number of hidden units selected from one to ten by validation-set mean
squared error. The nonlinear specification improves the in-sample fit by
$35\%$ relative to the linear SVAR.

The agent is a deep-RL controller that observes $x_t = (\pi_t, y_t, \cdot)$
and selects $i_t$ to maximise the discounted expected central-bank reward
$r_t = - w (\pi_t - \pi^\star)^2 - (1 - w) (y_t - y^\star)^2$. The policy is
parametrised either as a linear function of the state, in the spirit of a
Taylor rule, or as a nonlinear ANN. The estimated transition equations are
held fixed during policy optimisation, which is the assumption against which
the Lucas critique runs. The training algorithm is the deep-deterministic
policy gradient of \citet{LillicrapDDPG2016}, with the target networks and
replay buffer adapted to the macroeconometric setting. The data is quarterly
US data from 1987Q3 to 2007Q2, deliberately stopping before the global
financial crisis so that the zero-lower-bound period does not contaminate the
transition estimation.

The empirical results of \citet{HinterlangTaenzer2024} verify the abstract
claim of the paper. All RL-optimised reaction functions outperform the
realised federal funds rate path and the standard rule-based policies of
\citet{Taylor1993}, \citet{TaylorWilliams2010}, and the Federal Reserve's
Monetary Policy Report on the calibrated central-bank loss. Quantitatively,
the linear RL reaction functions reduce the central-bank loss by over $35\%$
relative to the realised rate, and the nonlinear RL reaction functions reduce
it by over $43\%$, with the additional eight percentage points coming from
the agent's ability to exploit the curvature in $f_y$ and $f_\pi$ that the
linear policy cannot match. The historical counterfactual path of inflation
and the output gap under the nonlinear RL rule sits visibly closer to the
$(\pi^\star, y^\star)$ targets than either the realised path or the standard
rules. A model-comparison exercise across eleven DSGE models confirms that
the linear RL rules remain stable under reasonable alternative
specifications, although the nonlinear rules are more sensitive to the choice
of structural model.

The Lucas critique applies in full force here. The transition equations are
estimated from a single historical sample with a particular monetary-policy
reaction function in place, and \citet{HinterlangTaenzer2024} explicitly
acknowledge that allowing rational-expectations agents to update their
expectations to the new RL-optimised policy would alter the transition
parameters. The DSGE-model-comparison exercise is the paper's best available
remedy, and the qualitative robustness of the linear RL rule across eleven
DSGE models is the operative defence. The deeper remedy, namely identifying
restrictions and instrumental variables for the structural shifts the agent
might induce, requires the kind of econometric identification machinery that
Section~\ref{sec:bridge} sketched and that
Section~\ref{sec:forecasting_rl_open} returns to. The
\citet{HinterlangTaenzer2024} result is therefore best read as a
counterfactual quantification of how much realised-loss reduction was on the
table during the 1987-2007 regime, holding the structural transition
equations fixed, rather than as a recommendation for current monetary
policy.

\subsubsection{Order execution and market microstructure}
\label{sec:app_execution}

The fourth application set is order execution and the related microstructure
problem of translating short-horizon directional forecasts into limit-order
trading rules. The forecasting input is either an exogenous high-frequency
signal of mid-quote returns or the agent's belief about how its own trades
move the price. The decision is a sequence of limit-order placements over a
discrete time grid. Both Micheli-Monod and the Nagy-Calliess-Zohren paper
(cited here as \citealp{Briola2023} for bibliography-key continuity, see the
note in the bib entry) work at the intersection of econometric price-impact
modelling and deep reinforcement learning.

\citet{MicheliMonod2024} formalise the optimal-execution problem with
transient price impact following \citet{Obizhaeva2013} and Alfonsi-Schied-Slynko
(2012). The agent's task is to liquidate an initial portfolio $X_0$ across
$N + 1$ trades on an equispaced grid, with the executed price
$P^\xi_t = P^0_t + \sum_{k \le t} \xi_{t_k} G(t - t_k)$ depending on past
trades through a decay kernel $G$. The expected profit
$V^\xi = -\mathbb E\bigl[\sum_k \xi_{t_k} \int_0^{\xi_{t_k}} G(0)\, ds\bigr]$
admits a closed-form optimum under \citet{Alfonsi2012} when the decay kernel
is strictly convex, namely
$\xi^\star = - X_0\, M^{-1} \mathbf 1 / (\mathbf 1^\top M^{-1} \mathbf 1)$
with $M_{ij} = G(|t_i - t_j|)$. The non-Markovian feature of the state, which
includes the entire history of past trades $\xi_{0:(t-1)}$, makes standard
deep RL slow to converge. \citet{MicheliMonod2024} address this with an
auxiliary $Q$-function defined on the projected state
$\hat\pi(s_t) = (t, X^\xi_t, \xi_{0:(t-1)})$ that drops the stochastic price
component and reduces the input space from four dimensions to three. Their
Proposition 3.1 establishes that the auxiliary $Q$-function satisfies an
auxiliary Bellman equation, and Proposition 3.2 establishes the corresponding
deterministic policy gradient, which lets the DDPG critic and actor exploit
the auxiliary representation. The empirical study trains the agent over
$30{,}000$ episodes per kernel and demonstrates convergence to the closed-form
$\xi^\star$ on exponential, power-law, and two linear-resilience kernels. The
online-learning experiment varies the kernel decay parameter $\rho$ from $1$
to $0.5$ or $1.5$ over $100{,}000$ episodes, and the executed-strategy reward
tracks the moving optimal within a small stable deviation attributable to
exploration noise.

The Nagy-Calliess-Zohren paper at \citet{Briola2023} addresses the
complementary problem of converting an exogenous high-frequency directional
signal into a limit-order trading rule. The agent operates in an OpenAI-Gym
limit-order-book environment built on the ABIDES market simulator that
replays NASDAQ price-time priority message data from the LOBSTER dataset.
The directional signal is a Dirichlet random variable
$d_t \in \Delta^2$ over the three classes
$\{$down, stationary, up$\}$, with the Dirichlet parameters depending on the
realised smoothed future return $r_{t+h}$ over a fixed horizon $h$ and the
noise level controlled by the parameter $a$. The state vector includes the
time remaining in the episode, the agent's cash balance and stock inventory,
the directional signal probabilities, and the level-one bid and ask prices
and volumes including the agent's own orders. The action space is
single-share limit-order placements at bid, mid-quote, or ask, plus a skip
action and an inventory-driven market order. The training algorithm is the
APEX-Dueling-DQN of \citet{Horgan2018Apex} with double Q-learning, dueling
networks, and asynchronous distributed prioritised experience replay. The
neural architecture combines three feed-forward layers with an LSTM layer
on a $100$-state context window for both the value and advantage streams.

The empirical case for the Nagy-Calliess-Zohren agent is on $4.5$ months of
AAPL training data (Jan-May 2012) and $1.5$ months of out-of-sample data
(May-June 2012), with a robustness check on $2022$ data. Three noise levels
are tested, namely $a = 1.1$ for high noise, $a = 1.3$ for medium noise, and
$a = 1.6$ for low noise. The RL agent dominates the heuristic baseline at all
three noise levels, with the differences in episodic mean return and in
Sharpe ratio statistically significant under a paired $t$-test at $p \ll 0.1$
across $31$ test episodes per noise level. The action distribution differs
across noise regimes. Under the noisiest signal ($a = 1.1$), the agent
skips $24.5\%$ of opportunities and places sells at the bid $95.4\%$ of the
time and buys at the ask $97.5\%$ of the time, namely aggressive marketable
limit orders. Under the cleanest signal ($a = 1.6$), the agent skips $7.8\%$
and uses passive mid-quote placements more frequently, trading off execution
probability against transaction cost. The robustness check on 2022 data
confirms that the strategy generalises across the decade-long regime shift,
with the high-noise case being the lone non-profitable scenario.

The methodological message across \citet{MicheliMonod2024} and the
Nagy-Calliess-Zohren paper is that the forecast and the execution are best
trained jointly even when the forecast is exogenous. In Micheli-Monod, the
forecast is the parametric model of price impact, and the closed-form
optimum acts as the oracle benchmark that the DDPG converges to without
prior knowledge of the kernel. In Nagy-Calliess-Zohren, the forecast is the
noisy directional signal, and the agent's contribution is to translate the
signal into a profitable order-placement rule that the rule-based heuristic
cannot. Both papers vindicate the chapter's recurring thesis. Decision-aware
training, namely either learning the policy end-to-end against the realised
cost or constructing an action-specific representation that exploits the
decision-induced structure of the state, outperforms predict-then-optimise
even when the forecast is held fixed.
